\documentclass[11pt]{article}

\usepackage[a4paper,margin=27mm]{geometry}
\usepackage{amsmath,amssymb,amsthm,mathtools}
\usepackage{graphicx}
\usepackage[T1]{fontenc}
\usepackage{lmodern}
\usepackage{microtype}
\usepackage{xcolor}
\usepackage[hidelinks]{hyperref}

\definecolor{prospectivered}{RGB}{185,0,0}
\newenvironment{prospective}{\color{prospectivered}}{}
\newcommand{\futurework}[1]{\textcolor{prospectivered}{#1}}
\newcommand{\C}{\mathbb C}
\newcommand{\Z}{\mathbb Z}
\newcommand{\T}{\mathbb T}
\newcommand{\norm}[1]{\left\lVert #1\right\rVert}
\newtheorem{theorem}{Theorem}[section]
\newtheorem{proposition}[theorem]{Proposition}
\newtheorem{corollary}[theorem]{Corollary}

\title{\textbf{From Spectral Hamiltonians to Tight Binding in Two and Three Dimensions:\\
Isotropic Locality on Orthogonal and Non-Orthogonal Lattices}}
\author{Eugene Joseph M. Ragasa\\
Department of Physics, De La Salle University, Manila, Philippines}
\date{Pre-results working manuscript}

\begin{document}
\maketitle

\begin{prospective}
\noindent\textbf{Working-color convention.} Red text identifies prospective work,
unexecuted calculations, open controls, conditional conclusions, or material that must
be replaced by qualified evidence. Red text is not a numerical result or scientific
claim.
\end{prospective}

\begin{abstract}
Band eigenvalues do not determine a matrix-valued tight-binding representation or its
real-space locality. We formulate a multidimensional validation sequence that begins
with a finite cell--orbital Hamiltonian, converts it unitarily to complete Bloch and
spectral representations, selects a qualified retained subspace and periodic frame, and
then constructs its complete Wannier/tight-binding representation. The construction
makes explicit the state-space identifications required before Hamiltonian subtraction.
A Laurent-polynomial obstruction proves that spectral projection does not preserve
finite hopping support in general. For the orthogonal flat-parent master family, a
factorized Bessel identity proves that the irreducible restricted-alignment residual is
exactly the normalized Fourier weight outside the onsite block in both two and three
dimensions. A shell-resolved Jacobi--Anger formula and factorial bound certify the
orthogonal Fourier tail without fitting a decay law. Square and simple-cubic rank-two
systems provide the corresponding isotropic baselines.
Non-orthogonal but physically isotropic extensions replace coordinate-radius shells by
metric-complete lattice shells and axis permutations by the applicable lattice point
group. Gauge-covariant equivalence, restricted global alignment, cyclic Fourier
conversion, interpolation, deliberate hopping truncation, and total consequences remain
separate channels.
\futurework{Numerical results and cross-lattice conclusions remain conditional on
independently verified, dimension-specific calculations.}
\end{abstract}

\section{Introduction}

Paper 1 establishes that pointwise band equivalence does not determine hopping locality
under a fixed identification of retained coordinates. Paper 2 asks how that distinction
extends across spatial dimension and lattice geometry while the retained rank remains
two. The purpose is not to infer a Hamiltonian from eigenvalues alone. It is to compare
complete Hamiltonian representations connected by declared isometries and gauges, and
to quantify the representation cost introduced by finite hopping support.

The gauge dependence and localization of Wannier representations are established
subjects, from analytic and exponential localization results
\cite{Kohn1959AnalyticWannier,Brouder2007ExponentialWannier} through practical
maximally localized constructions and Wannier90
\cite{Mostofi2008Wannier90,Marzari2012MaximallyLocalized}. Compact support is a
stronger algebraic property and need not coincide with maximal localization
\cite{Sathe2024CompactWannier}. Material-specific Wannier and strained-Hamiltonian
models likewise exist for graphene and other hexagonal systems
\cite{Jung2013GrapheneWannier,Fang2018Strained2D}, while environment-dependent
Hamiltonian mappings address geometry-responsive electronic structure
\cite{Zhang2022EquivariantHamiltonians}. The present work does not claim any of those
ingredients individually as new. Its narrower target is a controlled separation of
pointwise spectral equivalence, declared gauge alignment, metric-complete truncation,
and the physical hopping range needed at a fixed tolerance.

The study separates two questions:
\begin{enumerate}
 \item How are a finite cell--orbital operator, its Bloch representation, its complete
       eigensystem, a retained periodic frame, and its Wannier/TB blocks related?
 \item After those spaces have been aligned, how does a restricted representation
       choice alter the hopping support required at the same sampled spectral accuracy?
\end{enumerate}

\begin{prospective}
The intended evidence ladder begins with orthogonal square and simple-cubic systems and
then adds non-orthogonal but physically isotropic cells. Direction-dependent coupling
laws remain outside Paper 2. Retained rank greater than two belongs to Paper 3. None of
these stage assignments asserts that the corresponding calculation has been executed.
\end{prospective}

\section{Complete representation chain}

For dimension $d\in\{2,3\}$ and finite periodic size $N$, define
\begin{equation}
 \mathcal L_N^{(d)}=(\C^N)^{\otimes d},\qquad
 \mathcal H_{p,N}^{(d)}=\mathcal L_N^{(d)}\otimes\C^4,\qquad
 \mathcal H_{r,N}^{(d)}=\mathcal L_N^{(d)}\otimes\C^2.
 \label{eq:spaces}
\end{equation}
Kronecker translation operators act on $\mathcal L_N^{(d)}$; they do not tensor the
four-dimensional orbital fiber. With $\mathsf S_{\alpha,N}$ denoting translation on
axis $\alpha$, the orthogonal finite parent is
\begin{equation}
 \widehat K_{\rho,N}^{(d)}
 =I_{N^d}\otimes T_0+
 \frac{\rho}{d}\sum_{\alpha=1}^d
 \left(\mathsf S_{\alpha,N}\otimes T_+
 +\mathsf S_{\alpha,N}^\dagger\otimes T_-\right).
 \label{eq:kronecker-parent}
\end{equation}
The centered tensor Fourier transform $\mathcal F_{d,N}$ produces
\begin{equation}
 (\mathcal F_{d,N}\otimes I_4)\widehat K_{\rho,N}^{(d)}
 (\mathcal F_{d,N}^\dagger\otimes I_4)
 =\bigoplus_{\mathbf q}K_\rho^{(d)}(\mathbf q).
 \label{eq:bloch-map}
\end{equation}
For a complete fiber eigensystem
\begin{equation}
 K(\mathbf q)=Q(\mathbf q)E(\mathbf q)Q^\dagger(\mathbf q),
 \label{eq:eigensystem}
\end{equation}
$E(\mathbf q)$ is a diagonal spectral representation only together with the eigenvector
map $Q(\mathbf q)$. An eigenvalue array without its state-space map is not an operator in
orbital coordinates and cannot be subtracted from a TB Hamiltonian.

A qualified periodic frame $F(\mathbf q):\C^2\to\C^4$ for the retained projector defines
\begin{equation}
 H(\mathbf q)=F^\dagger(\mathbf q)K(\mathbf q)F(\mathbf q).
 \label{eq:retained}
\end{equation}
The inverse tensor Fourier transform of $\bigoplus_{\mathbf q}H(\mathbf q)$ gives the
complete retained Wannier/TB operator. Its block-circulant entries are the finite Fourier
hopping blocks $T_{\mathbf R}^{(N)}$.

\subsection{Spectral projection is not finite-range compression}

For a spectral projector $P_I$ with isometry $V_I$ onto its range,
\begin{equation}
 H_{\mathrm{red}}=V_I^\dagger H V_I,
 \qquad
 V_IH_{\mathrm{red}}V_I^\dagger=P_IHP_I.
 \label{eq:spectral-projection}
\end{equation}
This operation preserves the retained eigenpairs exactly, but it need not preserve a
finite hopping range. The following elementary example makes that statement exact.

\begin{proposition}[Spectral projection can destroy finite hopping support]
\label{prop:projection-destroys-range}
Let
\begin{equation}
 H(k)=\begin{pmatrix}
 \Delta&t(1+e^{-ik})\\
 t(1+e^{ik})&-\Delta
 \end{pmatrix},
 \qquad \Delta t\ne0.
 \label{eq:projection-example}
\end{equation}
Then $H$ has nearest-neighbor Fourier support and eigenvalues
$E_\pm(k)=\pm d(k)$, where
\begin{equation}
 d(k)=\sqrt{\Delta^2+4t^2\cos^2(k/2)}.
\end{equation}
The embedded lower-band operator
\begin{equation}
 P_-(k)H(k)P_-(k)=\tfrac12[H(k)-d(k)I]
 \label{eq:projected-range}
\end{equation}
is not a matrix-valued finite Fourier polynomial.
\end{proposition}

\begin{proof}
Set $z=e^{ik}$. If $d(k)$ were a finite Fourier polynomial, there would be a Laurent
polynomial $f(z)$ whose restriction to the unit circle equals $d(k)$. Squaring on the
unit circle and using equality on infinitely many points would give the Laurent-polynomial
identity
\begin{equation}
 f(z)^2=\Delta^2+2t^2+t^2(z+z^{-1}).
 \label{eq:laurent-square}
\end{equation}
If $m$ and $n$ are the smallest and largest exponents with nonzero coefficients in
$f$, then the smallest and largest exponents in $f^2$ are $2m$ and $2n$; their
coefficients cannot cancel because each is the square of an extreme coefficient. The
right-hand side of Eq.~\eqref{eq:laurent-square} has extreme exponents $-1$ and $1$
when $t\ne0$, which is impossible because $2m$ and $2n$ are even. Thus $d$ is not a
finite Fourier polynomial. Equation~\eqref{eq:projected-range} has diagonal entries
containing $-d(k)/2$ plus finite Fourier terms, so the embedded projected operator is not
finite range.
\end{proof}

The gap condition $\Delta\ne0$ keeps the two bands separated. The proposition does not
claim that every spectral projection of every finite-range parent is long ranged; it
shows that finite support is not preserved as an operation. Projection, frame selection,
Fourier conversion, hopping truncation, and restriction to a candidate model class are
therefore distinct operations.

\section{Orthogonal isotropic baselines}

The square and simple-cubic baselines use
\begin{equation}
 K_\rho^{(d)}(\mathbf q)=T_0+\frac{\rho}{d}
 \sum_{\alpha=1}^d
 \left(e^{2\pi iq_\alpha}T_+
 +e^{-2\pi iq_\alpha}T_-\right)
 \label{eq:orthogonal-parent}
\end{equation}
and
\begin{equation}
 U_a^{(d)}(\mathbf q)=R\!\left[
 \frac{a}{d}\sum_{\alpha=1}^d\sin(2\pi q_\alpha)
 \right].
 \label{eq:rotation}
\end{equation}
The flat-parent correlations are $J_0(a)^2$ in two dimensions and
$J_0(2a/3)^3$ in three dimensions. The next theorem proves these formulas and their
exact real-space meaning. They are continuum analytic comparators and do not replace
finite-mesh correlations. Square and cubic truncations retain every attainable metric
shell through the declared cap using complete signed-permutation orbits.

\subsection{Flat-parent residual--locality theorem}

Throughout, $E_{\mathrm{ref}}>0$ is an arbitrary reference-energy normalization, not a
band gap. Give $\T^d=[-1/2,1/2)^d$ normalized Lebesgue measure. Let
$D\in\mathbb R^{2\times2}$ be nonzero, symmetric, and traceless, and let the flat
retained parent be $H(\mathbf q)=\bar E I_2+D$. For the rotation in
Eq.~\eqref{eq:rotation}, define
\begin{align}
 B_a^{(d)}(\mathbf q)
 &=U_a^{(d)}(\mathbf q)^\dagger D U_a^{(d)}(\mathbf q),\\
 B_{\mathbf R}^{(d)}
 &=\int_{\T^d}e^{-2\pi i\mathbf q\cdot\mathbf R}
 B_a^{(d)}(\mathbf q)\,d\mathbf q,
 \qquad \mathbf R\in\Z^d,\\
 \kappa&=\frac{\norm{D}_F^2}{2E_{\mathrm{ref}}^2}.
 \label{eq:flat-definitions}
\end{align}
For one momentum-independent real rotation $R(\phi)$, centroid shift $s$, and the
nonnegative closure $\lambda\ge0$ of the splitting-scale domain, define
\begin{equation}
 L_{\mathrm{op},\phi}^2(s,\lambda)
 =\frac{1}{2E_{\mathrm{ref}}^2}\int_{\T^d}
 \norm{sE_{\mathrm{ref}}I_2+\lambda R(\phi)^\dagger D R(\phi)
       -B_a^{(d)}(\mathbf q)}_F^2\,d\mathbf q.
 \label{eq:flat-operator-loss}
\end{equation}
The normalized best global-rotation correlation is
\begin{equation}
 \mu_{\max}^{(d)}(a)=
 \max_{\phi\in\mathbb R}
 \frac{1}{\norm D_F^2}
 \int_{\T^d}
 \operatorname{tr}\!\left[
   R(\phi)^\dagger D R(\phi)B_a^{(d)}(\mathbf q)
 \right]d\mathbf q.
 \label{eq:flat-mu-definition}
\end{equation}

\begin{theorem}[Multidimensional flat-parent residual--locality identity]
\label{thm:flat-parent}
For $d\in\{2,3\}$ and $0\le a\le\pi/4$, set
\begin{equation}
 c_d(a)=J_0\!\left(\frac{2a}{d}\right)^d.
 \label{eq:flat-correlation-general}
\end{equation}
Then $c_d(a)>0$ and
\begin{align}
 \int_{\T^d}e^{2i\theta_a^{(d)}(\mathbf q)}\,d\mathbf q
 &=c_d(a),
 \label{eq:bessel-factorization}\\
 B_{\mathbf0}^{(d)}&=c_d(a)D.
 \label{eq:onsite-average}
\end{align}
The best global-rotation correlation is
$\mu_{\max}^{(d)}(a)=c_d(a)$, and the exact restricted-alignment residual satisfies
\begin{align}
 \inf_{s\in\mathbb R,\,\lambda\ge0,\,\phi\in\mathbb R}
 L_{\mathrm{op},\phi}^2(s,\lambda)
 &=\kappa\left[1-c_d(a)^2\right]
 \label{eq:flat-residual}\\
 &=\frac{1}{2E_{\mathrm{ref}}^2}
   \sum_{\mathbf R\ne\mathbf0}\norm{B_{\mathbf R}^{(d)}}_F^2.
 \label{eq:flat-tail}
\end{align}
If exact spectral agreement additionally fixes $s=0$ and $\lambda=1$, then
\begin{equation}
 \inf_{\phi\in\mathbb R}L_{\mathrm{op},\phi}^2(0,1)
 =2\kappa\left[1-c_d(a)\right].
 \label{eq:flat-exact-spectrum}
\end{equation}
Thus the two-dimensional formulas are
$c_2(a)=J_0(a)^2$ and the three-dimensional formulas are
$c_3(a)=J_0(2a/3)^3$.
\end{theorem}

\begin{proof}
A constant real rotation diagonalizes $D$ without changing any Frobenius norm, Fourier
weight, or minimization in the statement, so take $D=\delta\sigma_z$ with
$\delta\ne0$. Conjugation through a real rotation gives
\begin{equation}
 B_a^{(d)}(\mathbf q)
 =\delta\left[
 \cos\!\left(2\theta_a^{(d)}(\mathbf q)\right)\sigma_z
 -\sin\!\left(2\theta_a^{(d)}(\mathbf q)\right)\sigma_x
 \right].
 \label{eq:flat-rotated-pauli}
\end{equation}
Using the normalized integral representation
\begin{equation}
 J_0(x)=\int_{-1/2}^{1/2}e^{ix\sin(2\pi q)}\,dq
 \label{eq:bessel-integral}
\end{equation}
and the product structure of $\T^d$,
\begin{align}
 \int_{\T^d}e^{2i\theta_a^{(d)}(\mathbf q)}\,d\mathbf q
 &=\prod_{\alpha=1}^d
   \int_{-1/2}^{1/2}
   e^{i(2a/d)\sin(2\pi q_\alpha)}\,dq_\alpha\\
 &=J_0\!\left(\frac{2a}{d}\right)^d=c_d(a).
\end{align}
Global inversion $\mathbf q\mapsto-\mathbf q$ makes the sine average in
Eq.~\eqref{eq:flat-rotated-pauli} vanish, proving
Eq.~\eqref{eq:onsite-average}. Moreover,
$0\le2a/d\le\pi/4$, and the real part of the integrand in
Eq.~\eqref{eq:bessel-integral} is at least $\cos(2a/d)>0$; hence $c_d(a)>0$.

For any constant matrix $C$, the zero-mean property of
$B_a^{(d)}-B_{\mathbf0}^{(d)}$ gives the orthogonal decomposition
\begin{equation}
 \int_{\T^d}\norm{C-B_a^{(d)}(\mathbf q)}_F^2\,d\mathbf q
 =\int_{\T^d}\norm{B_a^{(d)}(\mathbf q)-B_{\mathbf0}^{(d)}}_F^2\,d\mathbf q
  +\norm{C-B_{\mathbf0}^{(d)}}_F^2.
 \label{eq:constant-projection}
\end{equation}
The constant $B_{\mathbf0}^{(d)}=c_d(a)D$ belongs to the allowed family in
Eq.~\eqref{eq:flat-operator-loss}: choose $s=0$, $\lambda=c_d(a)$, and $\phi=0$.
It is therefore the best constant restricted alignment. Equivalently, the normalized
maximum correlation with one global rotation is $c_d(a)$.

Parseval's identity applied entrywise yields
\begin{equation}
 \int_{\T^d}\norm{B_a^{(d)}-B_{\mathbf0}^{(d)}}_F^2
 =\sum_{\mathbf R\ne\mathbf0}\norm{B_{\mathbf R}^{(d)}}_F^2.
\end{equation}
Pointwise unitary invariance gives
$\norm{B_a^{(d)}(\mathbf q)}_F=\norm D_F$. A second use of Parseval, together with
Eq.~\eqref{eq:onsite-average}, therefore gives
\begin{equation}
 \sum_{\mathbf R\ne\mathbf0}\norm{B_{\mathbf R}^{(d)}}_F^2
 =\norm D_F^2-\norm{B_{\mathbf0}^{(d)}}_F^2
 =\norm D_F^2\left[1-c_d(a)^2\right].
\end{equation}
Dividing by $2E_{\mathrm{ref}}^2$ proves Eqs.~\eqref{eq:flat-residual} and
\eqref{eq:flat-tail}. Finally, fixing $s=0$ and $\lambda=1$ leaves the non-onsite
weight plus the onsite mismatch:
\begin{align}
 \inf_\phi L_{\mathrm{op},\phi}^2(0,1)
 &=\kappa[1-c_d(a)^2]+\kappa[1-c_d(a)]^2\\
 &=2\kappa[1-c_d(a)],
\end{align}
which proves Eq.~\eqref{eq:flat-exact-spectrum}.
\end{proof}

\begin{theorem}[Shell-resolved coefficients and factorial Fourier-tail bound]
\label{thm:orthogonal-fourier-tail}
Use the constant real rotation from the preceding proof so that
$D=\delta\sigma_z$, and set
\begin{equation}
 z=\frac{2a}{d},\qquad
 j_{\mathbf R}(z)=\prod_{\alpha=1}^d J_{R_\alpha}(z),\qquad
 p(\mathbf R)=\sum_{\alpha=1}^dR_\alpha\pmod 2.
\end{equation}
Then every Fourier block is known exactly:
\begin{equation}
 B_{\mathbf R}^{(d)}
 =\delta j_{\mathbf R}(z)
 \begin{cases}
  \sigma_z,&p(\mathbf R)=0,\\
  i\sigma_x,&p(\mathbf R)=1.
 \end{cases}
 \label{eq:orthogonal-fourier-block}
\end{equation}
For the complete squared-radius shell $s\in\mathbb N_0$, define
\begin{equation}
 W_{d,s}(z)=
 \sum_{\substack{\mathbf R\in\mathbb Z^d\\\norm{\mathbf R}_2^2=s}}
 \prod_{\alpha=1}^dJ_{R_\alpha}(z)^2.
 \label{eq:shell-weight}
\end{equation}
Its normalized hopping-block power and the total shell inventory satisfy
\begin{align}
 \frac{1}{2E_{\mathrm{ref}}^2}
 \sum_{\norm{\mathbf R}_2^2=s}\norm{B_{\mathbf R}^{(d)}}_F^2
 &=\kappa W_{d,s}(z),
 \label{eq:shell-power}\\
 \sum_{s=0}^{\infty}W_{d,s}(z)&=1.
 \label{eq:shell-partition}
\end{align}
For an integer shell cap $S\ge0$, let
\begin{equation}
 L=\left\lfloor\sqrt{\frac{S}{d}}\right\rfloor+1,
 \qquad x=\frac{|z|}{2}.
\end{equation}
Whenever $x<L+1$, the omitted normalized power has the rigorous bound
\begin{equation}
 \frac{1}{2E_{\mathrm{ref}}^2}
 \sum_{\norm{\mathbf R}_2^2>S}\norm{B_{\mathbf R}^{(d)}}_F^2
 \le
 \kappa\,
 \frac{2d\,e^{z^2/2}x^{2L}}
 {(L!)^2\left[1-x^2/(L+1)^2\right]}.
 \label{eq:orthogonal-factorial-tail}
\end{equation}
Thus the orthogonal flat-parent Fourier tail is certified directly from $a$, $d$, and
$S$; no fitted decay law is required.
\end{theorem}

\begin{proof}
The Jacobi--Anger identity gives
\begin{equation}
 e^{2i\theta_a^{(d)}(\mathbf q)}
 =\prod_{\alpha=1}^d
  \sum_{n\in\mathbb Z}J_n(z)e^{2\pi inq_\alpha}
 =\sum_{\mathbf R\in\mathbb Z^d}
  j_{\mathbf R}(z)e^{2\pi i\mathbf q\cdot\mathbf R}.
\end{equation}
Because $J_{-n}(z)=(-1)^nJ_n(z)$, the cosine coefficient vanishes for odd
$p(\mathbf R)$ and equals $j_{\mathbf R}(z)$ for even $p(\mathbf R)$; the sine
coefficient vanishes for even parity and equals $-ij_{\mathbf R}(z)$ for odd parity.
Substitution into Eq.~\eqref{eq:flat-rotated-pauli} proves
Eq.~\eqref{eq:orthogonal-fourier-block}. Since both Pauli matrices have squared
Frobenius norm two,
\begin{equation}
 \norm{B_{\mathbf R}^{(d)}}_F^2
 =\norm D_F^2\prod_{\alpha=1}^dJ_{R_\alpha}(z)^2.
 \label{eq:block-power-product}
\end{equation}
Equation~\eqref{eq:shell-power} follows by grouping a complete shell. Parseval applied
to the unit-modulus function $e^{iz\sin(2\pi q)}$ gives
$\sum_{n\in\mathbb Z}J_n(z)^2=1$; taking the $d$-fold product proves
Eq.~\eqref{eq:shell-partition}.

If $\norm{\mathbf R}_2^2>S$, at least one coordinate obeys $|R_\alpha|\ge L$.
The union bound and the one-dimensional Bessel partition therefore give
\begin{equation}
 \sum_{\norm{\mathbf R}_2^2>S}
 \prod_{\alpha=1}^dJ_{R_\alpha}(z)^2
 \le 2d\sum_{n=L}^{\infty}|J_n(z)|^2.
 \label{eq:tail-union}
\end{equation}
The defining Bessel series and $(n+m)!\ge n!m!$ imply, for $n\ge0$,
\begin{equation}
 |J_n(z)|
 \le \frac{x^n}{n!}\sum_{m=0}^{\infty}\frac{x^{2m}}{m!}
 =\frac{x^n}{n!}e^{x^2}.
\end{equation}
Moreover, $(L+m)!\ge L!(L+1)^m$, so when $x<L+1$,
\begin{equation}
 \sum_{n=L}^{\infty}\frac{x^{2n}}{(n!)^2}
 \le \frac{x^{2L}}{(L!)^2}
 \sum_{m=0}^{\infty}\left(\frac{x^2}{(L+1)^2}\right)^m
 =\frac{x^{2L}}
 {(L!)^2\left[1-x^2/(L+1)^2\right]}.
\end{equation}
Combining this estimate with Eq.~\eqref{eq:tail-union}, using
$e^{2x^2}=e^{z^2/2}$ and Eq.~\eqref{eq:block-power-product}, proves
Eq.~\eqref{eq:orthogonal-factorial-tail}.
\end{proof}

\begin{corollary}[Monotonicity and the small-amplitude dimension law]
\label{cor:dimension-law}
Define the minimum-alignment and exact-spectrum thresholds by
\begin{equation}
 \tau_{\min,d}^2(a)=\kappa[1-c_d(a)^2],
 \qquad
 \tau_{\mathrm{exact},d}^2(a)=2\kappa[1-c_d(a)].
\end{equation}
Both thresholds are strictly increasing for $0<a\le\pi/4$. As $a\to0$,
\begin{align}
 c_d(a)
 &=1-\frac{a^2}{d}+\frac{2d-1}{4d^3}a^4+O(a^6),
 \label{eq:correlation-small-a}\\
 \tau_{\min,d}^2(a)
 &=\frac{2\kappa}{d}a^2
   -\frac{\kappa(4d-1)}{2d^3}a^4+O(a^6),
 \label{eq:min-threshold-small-a}\\
 \tau_{\mathrm{exact},d}^2(a)
 &=\frac{2\kappa}{d}a^2
   -\frac{\kappa(2d-1)}{2d^3}a^4+O(a^6).
 \label{eq:exact-threshold-small-a}
\end{align}
In particular, both thresholds have leading slope
$\sqrt{2\kappa/d}$ with respect to the maximum rotation amplitude $a$.
\end{corollary}

\begin{proof}
For $x>0$, $J_0'(x)=-J_1(x)$. On $0<x\le\pi/4$, the alternating series
\begin{equation}
 J_1(x)=\sum_{m=0}^{\infty}
 \frac{(-1)^m}{m!(m+1)!}\left(\frac{x}{2}\right)^{2m+1}
\end{equation}
has strictly decreasing term magnitudes because their successive ratio is at most
$\pi^2/128<1$. Hence $J_1(x)>0$. Together with the positivity already proved for
$J_0(2a/d)$, this shows
\begin{equation}
 c_d'(a)=-2J_0(2a/d)^{d-1}J_1(2a/d)<0
\end{equation}
for $a>0$. The threshold monotonicity follows immediately. Finally,
$J_0(x)=1-x^2/4+x^4/64+O(x^6)$. Substituting $x=2a/d$, raising to the $d$th power, and
expanding the two threshold formulas proves
Eqs.~\eqref{eq:correlation-small-a}--\eqref{eq:exact-threshold-small-a}.
\end{proof}

Theorems~\ref{thm:flat-parent} and \ref{thm:orthogonal-fourier-tail}, together with
Corollary~\ref{cor:dimension-law}, apply to the orthogonal flat-parent master family.
The factorial bound certifies its complete-shell Fourier tail, but it must not be
transferred to a dependent non-orthogonal bond orbit. The bond-orbit theorem below
supplies that comparator's analytic form; a lattice-specific constrained-series tail
still requires a separate certificate.

\begin{prospective}
The orthogonal calculations belong to separate publication-independent 2-D and 3-D
experiment families. Their mesh sizes, shell caps, tolerances, convergence outcomes, and
evidence identities remain to be qualified independently and are not shared merely
because this manuscript compares them.
\end{prospective}

\subsection{Extracted two-dimensional precursor evidence}

An earlier retained scalar calculation used
\begin{equation}
 \widehat H_{2\mathrm D}=-\nabla^2+0.5\cos x+0.5\cos y
 +\lambda_{xy}\cos x\cos y.
 \label{eq:legacy-2d}
\end{equation}
At $\lambda_{xy}=0$, its represented Kronecker-sum discrepancy was zero and its
product-spectrum discrepancy was $3.55\times10^{-14}E_{\mathrm{ref}}$. At $\lambda_{xy}=0.30$, the
mixed-direction hopping norm was $1.94\times10^{-3}E_{\mathrm{ref}}$, compared with a
$3.00\times10^{-15}E_{\mathrm{ref}}$ numerical floor at zero coupling. On the strongest retained
coupling, the withheld RMS interpolation error decreased from
$5.75\times10^{-2}E_{\mathrm{ref}}$ for the onsite shell to $1.70\times10^{-6}E_{\mathrm{ref}}$ for the complete
finite transform. The latter reconstructed its own $15\times15$ mesh within
$1.11\times10^{-15}E_{\mathrm{ref}}$ while retaining nonzero error on an independent
$31\times31$ mesh.

These values are extracted from the preserved Paper 1 v2 controlled calculation. They
establish a historical finite-mesh separability, mixed-hopping, and off-node
interpolation control; they are not results of the new square-lattice rank-two protocol,
do not establish a 3-D outcome, and do not qualify a non-orthogonal lattice.

\section{Non-orthogonal but physically isotropic extension}

Let the direct-lattice basis be the columns of $A$ and define the direct-lattice metric
and reciprocal basis by
\begin{equation}
 G=A^{\mathsf T}A,
 \qquad B=2\pi A^{-\mathsf T}.
\end{equation}
For $\mathbf R\in\Z^d$, physical squared distance is
\begin{equation}
 r_G^2(\mathbf R)=\mathbf R^{\mathsf T}G\mathbf R.
 \label{eq:metric-radius}
\end{equation}
Thus the coordinate radius $R_x^2+R_y^2$ is not used for a non-orthogonal cell unless it
coincides with the declared metric. The integer lattice point group is
\begin{equation}
 \Gamma_G=\{M\in GL(d,\Z):M^{\mathsf T}GM=G\}.
 \label{eq:metric-group}
\end{equation}
An isotropic hopping inventory retains complete $\Gamma_G$ orbits and every attainable
metric shell through its cap. Equal coefficients are required only when imposed by the
model's point-group representation; support symmetry does not itself imply coefficient
equality.

A selected non-orthogonal parent must declare an inversion-complete and point-group-
complete bond orbit $\mathcal B$. A dimension-normalized candidate form is
\begin{equation}
 K_{\rho,G}(\mathbf q)=T_0+
 \frac{\rho}{|\mathcal B_+|}
 \sum_{\mathbf b\in\mathcal B_+}
 \left(e^{2\pi i\mathbf q\cdot\mathbf b}T_{\mathbf b}
 +e^{-2\pi i\mathbf q\cdot\mathbf b}T_{\mathbf b}^\dagger\right),
 \label{eq:general-parent}
\end{equation}
subject to the selected internal point-group representation. The smooth rotation must
likewise be constructed from a complete bond orbit rather than privileged coordinate
axes.

\subsection{Bond-orbit analytic comparator}

Let the chosen positive bond inventory be
$\mathcal B_+=\{\mathbf b_1,\ldots,\mathbf b_m\}\subset\Z^d\setminus\{\mathbf0\}$.
For real orbit weights $\eta_j$, set $W=\sum_j|\eta_j|>0$ and define the normalized
bond-orbit phase
\begin{equation}
 \theta_{a,\mathcal B}(\mathbf q)
 =\frac{a}{W}\sum_{j=1}^m\eta_j
   \sin(2\pi\mathbf q\cdot\mathbf b_j),
 \qquad |\theta_{a,\mathcal B}(\mathbf q)|\le a.
 \label{eq:bond-orbit-phase}
\end{equation}
The weights must be fixed by the selected symmetry contract; the theorem below does not
by itself declare them isotropic.

\begin{theorem}[Bond-orbit flat-parent comparator]
\label{thm:bond-orbit}
For $0\le a\le\pi/4$, define
\begin{equation}
 c_{\mathcal B,\boldsymbol\eta}(a)
 =\int_{\T^d}e^{2i\theta_{a,\mathcal B}(\mathbf q)}\,d\mathbf q.
 \label{eq:bond-orbit-correlation}
\end{equation}
Then $c_{\mathcal B,\boldsymbol\eta}(a)$ is real and strictly positive. If
\begin{align}
 D_{a,\mathcal B}(\mathbf q)
 &=R(\theta_{a,\mathcal B}(\mathbf q))^\dagger
   D R(\theta_{a,\mathcal B}(\mathbf q)),\\
 D_{\mathbf R}^{\mathcal B}
 &=\int_{\T^d}e^{-2\pi i\mathbf q\cdot\mathbf R}
   D_{a,\mathcal B}(\mathbf q)\,d\mathbf q,
\end{align}
then
\begin{align}
 D_{\mathbf0}^{\mathcal B}
 &=c_{\mathcal B,\boldsymbol\eta}(a)D,\\
 \inf_{s,\lambda\ge0,\phi}L_{\mathcal B,\phi}^2(s,\lambda)
 &=\kappa[1-c_{\mathcal B,\boldsymbol\eta}(a)^2]
 =\frac{1}{2E_{\mathrm{ref}}^2}\sum_{\mathbf R\ne\mathbf0}
   \norm{D_{\mathbf R}^{\mathcal B}}_F^2,
 \label{eq:bond-orbit-tail}\\
 \inf_\phi L_{\mathcal B,\phi}^2(0,1)
 &=2\kappa[1-c_{\mathcal B,\boldsymbol\eta}(a)],
\end{align}
where $L_{\mathcal B,\phi}$ is Eq.~\eqref{eq:flat-operator-loss} with
$B_a^{(d)}$ replaced by $D_{a,\mathcal B}$.

Moreover, the comparator has the exact zero-mode expansion
\begin{equation}
 c_{\mathcal B,\boldsymbol\eta}(a)
 =\sum_{\substack{\mathbf n\in\Z^m\\
                   \sum_{j=1}^m n_j\mathbf b_j=\mathbf0}}
   \prod_{j=1}^m
   J_{n_j}\!\left(\frac{2a\eta_j}{W}\right).
 \label{eq:bond-orbit-bessel-sum}
\end{equation}
If the bond vectors have no nontrivial integer relation, this reduces to
\begin{equation}
 c_{\mathcal B,\boldsymbol\eta}(a)
 =\prod_{j=1}^mJ_0\!\left(\frac{2a\eta_j}{W}\right).
 \label{eq:independent-bond-product}
\end{equation}
\end{theorem}

\begin{proof}
Equation~\eqref{eq:bond-orbit-phase} is odd under
$\mathbf q\mapsto-\mathbf q$, so the imaginary part of
Eq.~\eqref{eq:bond-orbit-correlation} integrates to zero. Since
$|2\theta_{a,\mathcal B}|\le\pi/2$, its cosine is nonnegative; it is positive in a
neighborhood of $\mathbf q=\mathbf0$, where the phase vanishes. Hence the integral is
strictly positive. Diagonalizing $D$ by one constant real rotation and repeating
Eq.~\eqref{eq:flat-rotated-pauli} gives
$D_{\mathbf0}^{\mathcal B}=c_{\mathcal B,\boldsymbol\eta}D$. The constant-projection
identity in Eq.~\eqref{eq:constant-projection} and entrywise Parseval then prove the two
residual formulas exactly as in Theorem~\ref{thm:flat-parent}.

For the explicit comparator, apply the Jacobi--Anger expansion to every bond harmonic:
\begin{equation}
 e^{i(2a\eta_j/W)\sin(2\pi\mathbf q\cdot\mathbf b_j)}
 =\sum_{n_j\in\Z}
 J_{n_j}\!\left(\frac{2a\eta_j}{W}\right)
 e^{2\pi i n_j\mathbf q\cdot\mathbf b_j}.
\end{equation}
The finite product of these absolutely convergent series may be integrated termwise.
Torus orthogonality retains exactly the tuples satisfying
$\sum_jn_j\mathbf b_j=\mathbf0$, which proves
Eq.~\eqref{eq:bond-orbit-bessel-sum}. If the integer kernel of the bond matrix is
trivial, only $\mathbf n=\mathbf0$ remains, giving
Eq.~\eqref{eq:independent-bond-product}.
\end{proof}

\begin{theorem}[Constrained-series complete-shell tail certificate]
\label{thm:constrained-tail}
Let
\begin{equation}
 g_{\mathbf R}=
 \sum_{\substack{\mathbf n\in\mathbb Z^m\\
                   \sum_jn_j\mathbf b_j=\mathbf R}}
 \prod_{j=1}^mJ_{n_j}\!\left(\frac{2a\eta_j}{W}\right)
 \label{eq:constrained-coefficient}
\end{equation}
be the scalar Fourier coefficient of $e^{2i\theta_{a,\mathcal B}}$. In the constant
basis where $D=\delta\sigma_z$,
\begin{equation}
 D_{\mathbf R}^{\mathcal B}
 =\frac{\delta}{2}\left[
 (g_{\mathbf R}+g_{-\mathbf R})\sigma_z
 +i(g_{\mathbf R}-g_{-\mathbf R})\sigma_x
 \right].
 \label{eq:constrained-matrix-block}
\end{equation}
Let $\mathcal S_S$ be every lattice displacement in the complete physical-metric shells
through $S$. If finite constrained sums $\widetilde g_{\mathbf R}$ obey certified bounds
$|g_{\mathbf R}-\widetilde g_{\mathbf R}|\le\epsilon_{\mathbf R}$ for every
$\mathbf R\in\mathcal S_S$, then
\begin{equation}
 \frac{1}{2E_{\mathrm{ref}}^2}
 \sum_{\mathbf R\notin\mathcal S_S}
 \norm{D_{\mathbf R}^{\mathcal B}}_F^2
 \le\kappa\left\{
 1-\sum_{\mathbf R\in\mathcal S_S}
 \left[\max\{0,|\widetilde g_{\mathbf R}|-\epsilon_{\mathbf R}\}\right]^2
 \right\}.
 \label{eq:constrained-tail-certificate}
\end{equation}

For the selected equal-weight triangular orbit
$\{(1,0),(0,1),(1,-1)\}$, with $z=2a/3$,
\begin{equation}
 g_{(r_1,r_2)}=
 \sum_{\ell\in\mathbb Z}
 J_{r_1-\ell}(z)J_{r_2+\ell}(z)J_\ell(z).
 \label{eq:triangular-coefficient}
\end{equation}
For the selected equal-weight FCC orbit
\begin{equation}
 \{(1,0,0),(0,1,0),(0,0,1),(-1,1,0),(-1,0,1),(0,-1,1)\},
\end{equation}
with $z=2a/6$,
\begin{align}
 g_{(r_1,r_2,r_3)}
 =\sum_{u,v,w\in\mathbb Z}
 &J_{r_1+u+v}(z)J_{r_2-u+w}(z)J_{r_3-v-w}(z)\notag\\
 &\times J_u(z)J_v(z)J_w(z).
 \label{eq:fcc-coefficient}
\end{align}
Set $x=|z|/2$ and
\begin{align}
 A_0(z)&=e^{x^2}(2e^x-1),\\
 A_K(z)&=
 \frac{2e^{x^2}x^{K+1}}
 {(K+1)![1-x/(K+2)]}.
 \label{eq:absolute-bessel-tail}
\end{align}
Truncating Eq.~\eqref{eq:triangular-coefficient} to $|\ell|\le K$ gives
$\epsilon_{\mathbf R}=A_K(z)$. Truncating Eq.~\eqref{eq:fcc-coefficient} to
$|u|,|v|,|w|\le K$ gives
$\epsilon_{\mathbf R}=3A_K(z)A_0(z)^2$.
\end{theorem}

\begin{proof}
Jacobi--Anger expansion and torus orthogonality give
Eq.~\eqref{eq:constrained-coefficient}. The symmetry
$e^{2i\theta(-\mathbf q)}=\overline{e^{2i\theta(\mathbf q)}}$ makes every
$g_{\mathbf R}$ real. Splitting the cosine and sine coefficients into
$(g_{\mathbf R}\pm g_{-\mathbf R})/2$ proves
Eq.~\eqref{eq:constrained-matrix-block}. Hence
\begin{equation}
 \frac{\norm{D_{\mathbf R}^{\mathcal B}}_F^2}{\norm D_F^2}
 =\frac{g_{\mathbf R}^2+g_{-\mathbf R}^2}{2}.
\end{equation}
Parseval for the unit-modulus scalar phase gives
$\sum_{\mathbf R}g_{\mathbf R}^2=1$. Because $\mathcal S_S$ is inversion complete,
the exact matrix tail normalized by $\norm D_F^2$ is
$1-\sum_{\mathbf R\in\mathcal S_S}g_{\mathbf R}^2$. The reverse triangle inequality
$|g_{\mathbf R}|\ge\max\{0,|\widetilde g_{\mathbf R}|-\epsilon_{\mathbf R}\}$
proves Eq.~\eqref{eq:constrained-tail-certificate}.

Solving the triangular constraint gives
$(n_1,n_2,n_3)=(r_1-\ell,r_2+\ell,\ell)$, which proves
Eq.~\eqref{eq:triangular-coefficient}. Taking
$(u,v,w)=(n_4,n_5,n_6)$ in the FCC constraint gives
\begin{equation}
 (n_1,n_2,n_3)=(r_1+u+v,r_2-u+w,r_3-v-w),
\end{equation}
which proves Eq.~\eqref{eq:fcc-coefficient}. Finally, the defining Bessel series gives
$|J_n(z)|\le e^{x^2}x^{|n|}/|n|!$. Summing it over all integers gives $A_0$; summing
over $|n|>K$ and bounding the remaining exponential-series tail geometrically gives
$A_K$. In the triangular remainder the other two Bessel factors have magnitude at most
one. In the FCC remainder, a union bound over the three free indices gives
$3A_KA_0^2$.
\end{proof}

For the selected maximum angle $a=\pi/4$, bounded binary64 evaluations with a declared
$10^{-13}$ per-coefficient numerical margin give normalized scalar-tail upper bounds
$2.8760\times10^{-9}$ for the triangular shell cap $S=21$ at $K=10$, and
$5.5358\times10^{-6}$ for the FCC shell cap $S=7$ at $K=9$. These values evaluate the
proved certificate but are not formal interval-arithmetic results, convergence evidence,
or material claims.

The metric and point group select the physically admissible complete bond orbit; the
integer relations among the selected reduced-coordinate bond vectors determine whether
the comparator factorizes. Thus the square and cubic products arise from independent
primitive-axis harmonics, whereas a bond orbit with integer dependencies generally has
the constrained Bessel sum in Eq.~\eqref{eq:bond-orbit-bessel-sum} instead.

\begin{prospective}
Separate non-orthogonal families now select a graphene-inspired triangular/honeycomb
2-D parent and a silicon-inspired FCC/diamond 3-D parent, together with matched Bravais-
control pairs at equal reciprocal-torus RMS gauge amplitude. ``Inspired'' denotes only
the declared Bravais geometry, two-site basis, and bipartite connectivity; neither parent
is a material parametrization. The dependent triangular and FCC control orbits give
respectively a one-index and a three-index constrained Bessel sum under
Theorems~\ref{thm:bond-orbit} and \ref{thm:constrained-tail}. Bounded Monkhorst--Pack
development actions and independent development verifiers now exist for all four
2-D/3-D families. The constrained-series certificate is evaluated at the development
caps, but qualified shell caps, full mesh ladders, control-parent gap certification,
formal interval evaluation, final tolerances and schemas, and frozen revisions remain
open. The square and cubic Bessel products must not be transferred unchanged to
these dependent bond orbits.
\end{prospective}

\begin{prospective}
A publication-independent extension now specifies frozen-gauge controls along a
unit-area rhombic triangular--square deformation and a unit-volume FCC--BCC Bain
path. The latter uses the crystallographic deformation associated with Bain's classic
martensite construction \cite{Bain1924Martensite}, but the proposed calculation freezes
the reduced-coordinate Hamiltonian and therefore is only a synthetic geometry stress
test. It must not be interpreted as a phase-stability, elastic, or material prediction.
Strain-responsive hopping laws, including the kinds of material-specific and
environment-dependent models represented by
Refs.~\cite{Fang2018Strained2D,Zhang2022EquivariantHamiltonians}, are separate
robustness controls rather than corrections to the frozen-gauge result.
\end{prospective}

\section{Aligned Hamiltonian subtraction}

For $\mathcal U=\bigoplus_{\mathbf q}U(\mathbf q)$, define
\begin{equation}
 \widehat A_B=\mathcal U^\dagger\widehat H_B\mathcal U.
\end{equation}
The corresponding Wannier-space unitary
\begin{equation}
 \mathcal G_W=(\mathcal F_{d,N}^\dagger\otimes I_2)
 \mathcal U(\mathcal F_{d,N}\otimes I_2)
\end{equation}
is generally nonlocal and satisfies
\begin{equation}
 \widehat A_W=\mathcal G_W^\dagger\widehat H_W\mathcal G_W.
 \label{eq:covariant}
\end{equation}
The covariantly aligned subtraction vanishes by construction. The restricted scientific
comparison replaces $\mathcal G_W$ by one cell-local global rotation
$I_{N^d}\otimes R(\phi)$. A parent-space comparison requires the retained projector and
an explicit lift. Every subtraction must state its source and target spaces, basis,
gauge, energy reference, units, geometry, and normalization.

\section{Error and evidence ledger}

The protocol reports separately:
\begin{enumerate}
 \item parent--retained reconstruction after projector lifting;
 \item covariant representation equivalence;
 \item restricted global-alignment residual;
 \item cyclic Nyquist conversion power;
 \item finite interpolation error;
 \item deliberate truncation error; and
 \item total reconstructed-Hamiltonian and spectral consequences.
\end{enumerate}
Writing $H^{\mathrm{dir}}$ for the direct withheld-mesh field,
$H_N^{\mathrm{ord}}$ for the full ordinary Fourier interpolant, and
$H_{N,R}^{\mathrm{ord}}$ for its complete-radius truncation, the three approximation
channels are
\begin{align}
 \mathcal E_{\mathrm{interp}}&=\mathcal E(H_N^{\mathrm{ord}},H^{\mathrm{dir}}),\\
 \mathcal E_{\mathrm{trunc}}&=\mathcal E(H_{N,R}^{\mathrm{ord}},H_N^{\mathrm{ord}}),\\
 \mathcal E_{\mathrm{total}}&=\mathcal E(H_{N,R}^{\mathrm{ord}},H^{\mathrm{dir}}),
 \label{eq:three-error-channels}
\end{align}
where $\mathcal E$ is reported independently as maximum/RMS normalized Frobenius error
and maximum/RMS sorted-band error. These channels are not assumed additive:
interpolation and truncation may reinforce or cancel. Omitted hopping-block power is a
locality diagnostic, not a universal spectral-error bound. First-passing support
decisions are sampled decisions and use hatted notation. For a future strain-responsive
perturbation, Eq.~\eqref{eq:three-error-channels} must be applied again to the
 gauge-aligned difference field $\Delta H_s$, after constructing that field in one
common state space. The M0 strain plots report Fourier-tail locality only and do not
constitute a three-channel strain-potential calculation; the separately labeled M1/M2
smoke controls repeat the three channels without upgrading their evidence status.

Calculation parity with Paper~1 requires every applicable predecessor diagnostic, not
only the three error channels. In particular, the numerical alignment must report the
splitting scale and finite-mesh maximum correlation, minimize over one global rotation
\emph{and} the positive splitting scale for the minimum operator threshold, and report
the distinct unit-splitting exact-spectrum joint threshold. A fixed-unit-scale rotation
fit cannot substitute for those two Paper~1 quantities. The complete predecessor
catalog also requires frame/sewing, locality-power, Weyl, first-passing-context,
refinement, and provenance counterparts in every family.

\section{Results}

\begin{prospective}
Replaceable bounded development records exercise a partial numerical channel sequence
from square lattice through the silicon-inspired diamond target. A strict Paper~1
catalog audit initially found missing or incompletely serialized predecessor channels.
The regenerated records now map all 110 predecessor quantities as implemented,
explicitly unresolved, explicitly inapplicable, or multidimensional-partial, including
the corrected positive-splitting alignment and analytic amplitude-resolution ladder.
Row-level parity verification and scalar closure-phase aggregation remain open and still
preclude quantitative scientific use. The sparse actions use three
construction levels: $16^2$, $32^2$, and $64^2$ in 2-D and $8^3$, $16^3$, and $32^3$
in 3-D, with separately shifted evaluation meshes and unchanged models, thresholds, and
shell caps. All development verifiers pass with stored identity defects below
$2.0\times10^{-14}$; this envelope now includes the newly serialized transport and
closure-unitarity channels. One of the 18 graphene-family pointwise first-passing decision sets
moves by one shell between the first two levels, while every decision is unchanged on
the final refinement in all four families. Several strict decisions remain unresolved
at the deliberately small shell caps. These records establish software and bookkeeping
behavior only; row-level Paper~1 parity verification, frozen tolerances, and cap policy
remain open, so the three mesh levels do not yet populate the paper's qualified scientific
results. Figures~\ref{fig:family-square}--
\ref{fig:family-silicon-inspired} show the four distinct model-family development
summaries. They must not be read as four strain-responsive models.

\begin{figure}[p]
 \centering
 \includegraphics[width=0.98\textwidth]{../../../../../experiments/model-adequacy/periodic2d/band-equivalence-versus-hopping-locality/analysis/band-equivalence-locality-scientific-story.pdf}
 \caption{Replaceable square-lattice family development summary. The panels keep
 withheld interpolation, external retained-subspace isolation, hopping-tail power,
 withheld spectral consequence, and tail-power/consequence comparison separate.}
 \label{fig:family-square}
\end{figure}

\begin{figure}[p]
 \centering
 \includegraphics[width=0.98\textwidth]{../../../../../experiments/model-adequacy/periodic2d_nonorthog/band-equivalence-versus-hopping-locality/analysis/band-equivalence-locality-scientific-story.pdf}
 \caption{Replaceable graphene-inspired family development summary, including its
 matched Bravais controls. ``Inspired'' denotes geometry and connectivity only; this is
 not a graphene parametrization or validation result.}
 \label{fig:family-graphene-inspired}
\end{figure}

\begin{figure}[p]
 \centering
 \includegraphics[width=0.98\textwidth]{../../../../../experiments/model-adequacy/periodic3d/band-equivalence-versus-hopping-locality/analysis/band-equivalence-locality-scientific-story.pdf}
 \caption{Replaceable simple-cubic family development summary. The two displayed mesh
 levels establish software behavior only and are not a convergence result.}
 \label{fig:family-simple-cubic}
\end{figure}

\begin{figure}[p]
 \centering
 \includegraphics[width=0.98\textwidth]{../../../../../experiments/model-adequacy/periodic3d_nonorthog/band-equivalence-versus-hopping-locality/analysis/band-equivalence-locality-scientific-story.pdf}
 \caption{Replaceable silicon-inspired family development summary, including its
 matched Bravais controls. ``Inspired'' denotes geometry and connectivity only; this is
 not a silicon parametrization, material calculation, or Paper~4 substitute.}
 \label{fig:family-silicon-inspired}
\end{figure}

Figures~\ref{fig:three-channel-square}--\ref{fig:three-channel-silicon-inspired}
show the corresponding refined-mesh three-channel decompositions. Near overlap of the
deliberate-truncation and total curves means interpolation is much smaller at these
sampled points; it does not make the channels definitionally identical.

\begin{figure}[p]
 \centering
 \includegraphics[width=0.98\textwidth]{../../../../../experiments/model-adequacy/periodic2d/band-equivalence-versus-hopping-locality/analysis/three-channel-error-decomposition.pdf}
 \caption{Square-lattice endpoint comparison: maximum Hamiltonian and sorted-band
 interpolation, deliberate-truncation, and total error channels.}
 \label{fig:three-channel-square}
\end{figure}

\begin{figure}[p]
 \centering
 \includegraphics[width=0.98\textwidth]{../../../../../experiments/model-adequacy/periodic2d_nonorthog/band-equivalence-versus-hopping-locality/analysis/three-channel-error-decomposition.pdf}
 \caption{Graphene-inspired family and matched controls: refined-mesh endpoint
 three-channel decomposition. Every row retains its own complete physical-metric shell
 inventory.}
 \label{fig:three-channel-graphene-inspired}
\end{figure}

\begin{figure}[p]
 \centering
 \includegraphics[width=0.98\textwidth]{../../../../../experiments/model-adequacy/periodic3d/band-equivalence-versus-hopping-locality/analysis/three-channel-error-decomposition.pdf}
 \caption{Simple-cubic endpoint comparison: maximum Hamiltonian and sorted-band
 interpolation, deliberate-truncation, and total error channels.}
 \label{fig:three-channel-simple-cubic}
\end{figure}

\begin{figure}[p]
 \centering
 \includegraphics[width=0.98\textwidth]{../../../../../experiments/model-adequacy/periodic3d_nonorthog/band-equivalence-versus-hopping-locality/analysis/three-channel-error-decomposition.pdf}
 \caption{Silicon-inspired family and matched controls: refined-mesh endpoint
 three-channel decomposition. These are synthetic development channels, not silicon
 material errors.}
 \label{fig:three-channel-silicon-inspired}
\end{figure}

\paragraph{Bounded M0 observation.}
Separate frozen-parent geometry-stress diagnostics use nine unit-area rhombic and
unit-volume Bain metrics. The reduced-coordinate Hamiltonian, gauge, Fourier-block
weights, and restricted-alignment residual are invariant by construction. Nevertheless,
the changing physical metric reorders complete hopping shells and changes the sampled
first-passing physical radius and retained block count. At tail-power tolerance
$10^{-5}$, the endpoint radius changes from $3.7224$ to $4.2426$ in 2-D and from
$2.9698$ to $3.2734$ in 3-D; the corresponding block counts change from $43$ to $61$
and from $135$ to $169$. At the looser 3-D tolerance $10^{-3}$, the endpoint radius
instead changes from $2.2449$ to $2.1822$ while the block count changes from $55$ to
$59$. Physical range and parameter count are therefore distinct locality costs even
when the abstract reduced-coordinate operator is unchanged. A unimodular map identifies
the existing FCC primitive coordinates with the Bain endpoint coordinates.
Figures~\ref{fig:m0-strain-2d} and \ref{fig:m0-strain-3d} display the complete bounded
M0 diagnostics.

\begin{figure}[p]
 \centering
 \includegraphics[width=0.98\textwidth]{../../../../../experiments/model-adequacy/periodic2d/gauge-locality-under-synthetic-strain/analysis/development/gauge-locality-under-synthetic-strain.pdf}
 \caption{Replaceable 2-D frozen-gauge M0 development diagnostic. The parent, gauge,
 Fourier weights, and unit cell area remain fixed while the rhombic metric moves from
 triangular to square. Panels separate bond-orbit splitting, invariant restricted
 residual, first-passing physical radius, and retained block count. The plotted
 tail-power thresholds are locality diagnostics, not universal spectral-error bounds.}
 \label{fig:m0-strain-2d}
\end{figure}

\begin{figure}[p]
 \centering
 \includegraphics[width=0.98\textwidth]{../../../../../experiments/model-adequacy/periodic3d/gauge-locality-under-synthetic-strain/analysis/development/gauge-locality-under-synthetic-strain.pdf}
 \caption{Replaceable 3-D frozen-gauge M0 development diagnostic along the unit-volume
 FCC--BCC Bain metric path. The looser threshold shows that first-passing physical
 radius and retained block count need not move together. This synthetic geometry stress
 test is not a phase-stability, elasticity, strain-potential, or material result.}
 \label{fig:m0-strain-3d}
\end{figure}

\paragraph{Bounded M1/M2 response controls.}
A separately labeled smoke action applies the sampled-qualified four-state bipartite parent.
M1 freezes its onsite and hopping matrices; its pointwise gauge-aligned strain
difference remains below $1.4\times10^{-15}E_{\mathrm{ref}}$ on both paths. M2 changes
only the parent-bond amplitudes through the prospectively fixed dimensionless radial law
$t_j(s)/t_j(0)=\exp[-(d_j(s)/d_j(0)-1)]$. At the final 2-D metric, the responses span
$0.81758$--$1.27146$ and the maximum aligned retained-space perturbation is
$0.063575E_{\mathrm{ref}}$. At the BCC endpoint, the four selected bonds share response
$0.97564$ and the corresponding maximum is $0.007289E_{\mathrm{ref}}$. The stored
trace/traceless decomposition is dominated by the traceless channel in both smoke
records. At the $10^{-5}$ tail support, the 2-D endpoint interpolation, deliberate-
truncation, and total maximum matrix errors of $\Delta H_s$ are approximately
$6.0\times10^{-8}$, $7.801\times10^{-4}$, and $7.802\times10^{-4}$; the 3-D values
are $9.36\times10^{-5}$, $1.61\times10^{-4}$, and $1.89\times10^{-4}$. These
channels are displayed in Figs.~\ref{fig:m1m2-strain-2d} and
\ref{fig:m1m2-strain-3d} and are not added.

\begin{figure}[p]
 \centering
 \includegraphics[width=0.98\textwidth]{../../../../../experiments/model-adequacy/periodic2d/gauge-locality-under-synthetic-strain/analysis/response-development/m1-m2-strain-response.pdf}
 \caption{Replaceable 2-D M1/M2 response smoke diagnostic. M1 verifies the zero aligned
 difference expected from the frozen parent. M2 uses the fixed synthetic radial law and
 reports its aligned perturbation, scalar/traceless split, locality cost, and three
 approximation channels.}
 \label{fig:m1m2-strain-2d}
\end{figure}

\begin{figure}[p]
 \centering
 \includegraphics[width=0.98\textwidth]{../../../../../experiments/model-adequacy/periodic3d/gauge-locality-under-synthetic-strain/analysis/response-development/m1-m2-strain-response.pdf}
 \caption{Replaceable 3-D M1/M2 Bain-response smoke diagnostic. The fixed M2 radial law
 is a synthetic sensitivity control, not a phase-stability or material-silicon model.}
 \label{fig:m1m2-strain-3d}
\end{figure}

These single-transform M0 and single-smoke M1/M2 records lack retained schemas,
independent artifact verification, and convergence and therefore remain software-
development diagnostics, not a material strain potential, material prediction, or paper
evidence. The extracted scalar
2-D precursor retains its historical finite-mesh scope. New scientific claims require
qualified calculations from the cited experiment families.
\end{prospective}

\section{Bounded interpretation}

\begin{prospective}
If supported by future evidence, the paper may compare representation cost across the
selected orthogonal and non-orthogonal isotropic systems. It will not infer a universal
range law, claim material validity, identify eigenvalue equality with operator equality,
or transfer a cutoff or tolerance between lattice geometries without qualification.
\end{prospective}

\bibliographystyle{unsrt}
\bibliography{references}

\end{document}
